\documentclass[11pt]{article}
\usepackage[margin=1in]{geometry}
\usepackage{amsmath,amssymb}
\usepackage{booktabs}
\usepackage{graphicx}
\usepackage{xcolor}
\usepackage{hyperref}

% Placeholder macro for numbers/text still to be filled in (e.g. after the seeded fine-tune run).
\newcommand{\pending}[1]{\textcolor{red}{[#1]}}
\newcommand{\todo}[1]{\textcolor{orange}{\textbf{TODO:} #1}}

\title{Do MEG Foundation Models Already Encode Alzheimer's Disease?\\
\large BrainOmni evaluation (draft --- BrainOmni portion)}
\author{Fae \and Srikantan Nagarajan \\ \small (MEG-GPT portion by Sid, to be integrated)}
\date{Draft \today}

\begin{document}
\maketitle

\begin{abstract}
\todo{Write last, after Results are final and the MEG-GPT half is integrated.}
We evaluate whether a self-supervised MEG foundation model, BrainOmni, implicitly encodes
representations that distinguish Alzheimer's disease (AD) from cognitively healthy controls,
without any disease-specific pretraining. On a 148-subject resting-state MEG cohort, frozen
BrainOmni features classified AD at 0.831 accuracy (AUC 0.87), and task-specific fine-tuning of the
last 1--3 transformer blocks did not improve on this (0.79--0.82 accuracy). Both were statistically
indistinguishable, on overlapping 95\% bootstrap confidence intervals, from a simple six-feature
handcrafted spectral/connectivity baseline (0.872 accuracy, AUC 0.93). Fine-tuning learning curves
show rapid overfitting that worsens with depth, and BrainOmni remains on the rising portion of its
sample-size learning curve. These results indicate that, on this cohort, the foundation model does
not capture AD-discriminative structure beyond established domain features, and appears
data/representation-limited rather than better or worse than handcrafted analysis.
\end{abstract}

\section{Introduction}
\todo{Draft after Methods/Results. Key beats: (1) foundation models promise transfer with little
downstream data; (2) MEG + AD context and the clinical stakes; (3) the open question --- do these
models already encode disease-relevant structure without disease-specific training?; (4) our
contribution: a careful, reproducible negative result on a real AD cohort, evaluating frozen and
fine-tuned BrainOmni against a strong handcrafted baseline with proper uncertainty quantification.}

\section{Methods}

\subsection{Dataset and preprocessing}
We use the resting-state MEG cohort of Kudo et al.\ \cite{kudo2023meg}, collected at UCSF:
78 patients with Alzheimer's disease and 70 cognitively healthy controls ($N=148$), recorded on a
whole-head gradiometer system ($\sim$271 channels) at 600\,Hz.

Each recording was resampled to 256\,Hz using polyphase filtering (which applies an anti-aliasing
low-pass), and each channel was $z$-normalized to zero mean and unit variance over the recording.
Sensor positions and orientations were read from the gradiometer geometry to form a $(C,6)$ array
per subject, mean-centered and scale-normalized; all channels were typed as gradiometers to match
BrainOmni's sensor encoder. From each recording we extracted two non-overlapping 30-second segments
(60\,s total), which are tokenized separately and averaged at the feature level (below). We applied
no additional band-pass or notch filtering.
\todo{Confirm with Sri whether the eLife recordings were already band-pass/notch filtered; BrainOmni's
own preprocessing applies a 0.1--96\,Hz band-pass and 50/60\,Hz notch, which we did not repeat.}

\subsection{Foundation model: BrainOmni}
BrainOmni is a self-supervised transformer for neural time series, pretrained on
$\sim$1{,}997\,h of EEG and $\sim$656\,h of MEG, none of it AD-specific. A learned BrainTokenizer
maps arbitrary sensor montages to a fixed set of 16 ``neuro tokens'' per time window (via Perceiver
queries), so the representation dimensionality is independent of the $\sim$271 input MEG channels.
We evaluate two pretrained sizes with identical tokenizers and depth (12 transformer blocks):
\textbf{tiny} (model width $256$) and \textbf{base} (model width $512$).

\subsection{Frozen feature extraction}
For frozen evaluation we reproduce the model's \texttt{encode} path: tokenize, add the learned
neuro embeddings, project, pass through the transformer blocks, and $L_2$-normalize the final
representation. Per subject we average over time windows and over the two segments, then flatten the
16 tokens, yielding a feature vector of dimension $16\times256=4096$ (tiny) or $16\times512=8192$
(base). These are the frozen features fed to a linear classifier.

\subsection{Fine-tuning}
For fine-tuning we unfreeze the last 1, 2, or 3 transformer blocks and attach a small classification
head (Linear $\rightarrow$ ReLU $\rightarrow$ Dropout(0.3) $\rightarrow$ Linear). Following the
model's official downstream recipe we use differential learning rates (pretrained backbone
$1\times10^{-5}$, randomly initialized head $1\times10^{-4}$), AdamW (weight decay $0.05$,
$\beta=(0.9,0.99)$), a warmup-cosine schedule (10\% warmup), label smoothing $0.1$, and gradient
clipping at $1.0$, trained for 200 epochs with the reported weights taken from the epoch of lowest
inner-validation loss.

\subsection{Classification heads}
The fine-tuned backbone is read out with two heads on identical features for a fair comparison:
the jointly trained neural-network (NN) head, and a logistic-regression (LR) head fit on the
extracted fine-tuned features. Frozen features are classified with LR.

\subsection{Handcrafted baseline}
As a domain baseline we use six features previously reported for this cohort: amplitude envelope
correlation (AEC) and spectral power, each in the delta--theta, alpha, and beta bands, expressed as
control-normative $z$-scores. We verified this baseline is robust to normalization: recomputing the
$z$-scores within each training fold (no test-set leakage) reproduces the same accuracy.

\subsection{Evaluation protocol}
All methods are evaluated with nested stratified 5-fold cross-validation using a fixed seed (123),
so every method sees identical folds. Within each training fold a 15\% inner-validation split is
carved out for epoch selection; the held-out test fold is never used for selection (it is monitored
only for diagnostic learning curves). We report accuracy and area under the ROC curve (AUC).
Uncertainty is quantified by bootstrapping the 148 pooled out-of-fold predictions
(2{,}000 resamples; 95\% percentile intervals). Network training is seeded (123) for reproducibility.
Robustness of the linear-classifier choice is assessed with (i) a sweep of $\sim$23 classifiers and
(ii) a PCA-dimensionality sweep (1--50 components, fit within each fold).

\section{Results}

\subsection{Frozen BrainOmni does not beat a handcrafted baseline}
Table~\ref{tab:frozen} reports 5-fold CV performance. Frozen BrainOmni reached 0.831 accuracy for
both model sizes, below the handcrafted baseline's 0.872. Interestingly, the larger \emph{base}
model gave no accuracy gain over \emph{tiny} but a modest AUC gain (0.878 vs.\ 0.866), indicating the
wider model improves ranking slightly without changing the operating-point accuracy.

\begin{table}[h]
\centering
\begin{tabular}{|l|c|c|}
\hline
Method & Accuracy [95\% CI] & AUC [95\% CI] \\
\hline
\multicolumn{3}{|l|}{\textit{Handcrafted}} \\
\hline
\quad 6 features + LR              & \textbf{0.872} [0.817, 0.926] & \textbf{0.925} [0.876, 0.968] \\
\hline
\multicolumn{3}{|l|}{\textit{Frozen BrainOmni}} \\
\hline
\quad base + LR                    & 0.831 [0.764, 0.885] & 0.878 [0.817, 0.929] \\
\hline
\quad tiny + LR                    & 0.831 [0.770, 0.892] & 0.866 [0.802, 0.921] \\
\hline
\end{tabular}
\caption{Frozen BrainOmni vs.\ the handcrafted baseline, 5-fold CV (seed 123, $N=148$), with
95\% bootstrap confidence intervals (2{,}000 resamples of the pooled out-of-fold predictions).
Confidence intervals overlap: frozen BrainOmni is not statistically distinguishable from
handcrafted.}
\label{tab:frozen}
\end{table}

\subsection{Fine-tuning does not help}
Table~\ref{tab:finetune} reports the fine-tuned (last-1-block) results alongside the frozen and
handcrafted numbers for reference. Fine-tuning the last 1--3 transformer blocks did \emph{not}
improve over the frozen baseline. The best fine-tuned configuration (base, last-1 block) reached
0.818 accuracy, still below both frozen (0.831) and handcrafted (0.872); the tiny model fine-tuned
to only 0.791--0.797. Fine-tuning \emph{more} blocks did not help and if anything hurt slightly:
across both models the ordering was monotone or flat in trainable depth (last-1 $\geq$ last-2
$\geq$ last-3; e.g.\ tiny 0.790, 0.783, 0.777). Fine-tuning also did not degrade ranking (fine-tuned
AUC 0.86--0.88, comparable to frozen), so the loss is at the operating point rather than in
separability. These results are consistent with the limited number of training subjects (see
Section~\ref{sec:datalimited}).

\begin{table}[h]
\centering
\begin{tabular}{|l|c|c|}
\hline
Method & Accuracy [95\% CI] & AUC [95\% CI] \\
\hline
\multicolumn{3}{|l|}{\textit{Reference (from Table~\ref{tab:frozen})}} \\
\hline
\quad \textit{Handcrafted + LR}     & \textit{0.872} [0.817, 0.926] & \textit{0.925} [0.876, 0.968] \\
\hline
\quad \textit{Frozen base + LR}     & \textit{0.831} [0.764, 0.885] & \textit{0.878} [0.817, 0.929] \\
\hline
\multicolumn{3}{|l|}{\textit{Fine-tuned BrainOmni (last-1 block)}} \\
\hline
\quad base + NN                    & 0.818 [0.757, 0.878] & 0.880 [0.819, 0.934] \\
\hline
\quad base + LR                    & 0.818 [0.750, 0.878] & 0.877 [0.815, 0.932] \\
\hline
\quad tiny + LR                    & 0.797 [0.736, 0.858] & 0.872 [0.808, 0.928] \\
\hline
\quad tiny + NN                    & 0.791 [0.723, 0.858] & 0.860 [0.797, 0.917] \\
\hline
\end{tabular}
\caption{Fine-tuned BrainOmni (last-1 block), 5-fold CV (seed 123, $N=148$), with 95\% bootstrap
confidence intervals. Reference rows (italicized) repeat the frozen and handcrafted numbers from
Table~\ref{tab:frozen} for comparison. All confidence intervals overlap.}
\label{tab:finetune}
\end{table}

Table~\ref{tab:depth} reports every fine-tuning depth. Accuracy is non-increasing with the number
of trainable blocks in all four configurations, i.e.\ unfreezing more of the backbone never helped
and in most cases hurt.

\begin{table}[h]
\centering
\begin{tabular}{|l|c|c|c|}
\hline
Configuration & Last-1 & Last-2 & Last-3 \\
\hline
tiny + NN & 0.790 & 0.783 & 0.777 \\
\hline
tiny + LR & 0.797 & 0.790 & 0.777 \\
\hline
base + NN & 0.817 & 0.817 & 0.783 \\
\hline
base + LR & 0.817 & 0.817 & 0.790 \\
\hline
\end{tabular}
\caption{Fine-tuning depth ablation: 5-fold CV accuracy (mean across folds) as a function of the
number of unfrozen transformer blocks, seed 123, $N=148$. Confidence intervals are omitted here
because the comparison of interest is the trend across depth rather than per-cell significance;
see Table~\ref{tab:finetune} for intervals on the last-1 configuration.}
\label{tab:depth}
\end{table}

\subsection{All methods are statistically indistinguishable}
Bootstrap 95\% confidence intervals overlap across all seven methods (Figure~\ref{fig:bootstrap};
Tables~\ref{tab:frozen} and \ref{tab:finetune}). Overlapping intervals are, however, only weak
evidence of equivalence: two 95\% intervals can overlap while the difference between the methods is
still significant. We therefore test each method directly against the handcrafted baseline with a
\emph{paired} bootstrap, resampling subjects once per iteration and evaluating both methods on the
same resample, so that the correlation induced by evaluating every method on identical subjects is
accounted for.

Table~\ref{tab:diff} reports these differences, using the last-1-block fine-tuned models: this is
the best-performing depth (Table~\ref{tab:depth}), so the comparison is deliberately the most
favorable one for fine-tuning, and the deeper configurations sit further below handcrafted still.
Every interval contains zero, for both accuracy and AUC, so no BrainOmni configuration, frozen or
fine-tuned, is statistically distinguishable from six handcrafted spectral/connectivity features.
The wide intervals are a direct consequence of the small
cohort ($N=148$): the data can resolve that all methods lie in the same $\sim$0.79--0.87 band, but
not finer differences between them. The honest conclusion is therefore not that any method ``wins,''
but that a general-purpose MEG foundation model is \emph{statistically tied with} an established
domain baseline on this task.

\begin{table}[h]
\centering
\begin{tabular}{|l|c|c|}
\hline
Method & $\Delta$ Accuracy vs handcrafted & $\Delta$ AUC vs handcrafted \\
\hline
Frozen base + LR      & \pending{$-$0.000 [$-$0.000, $+$0.000]} & \pending{---} \\
\hline
Frozen tiny + LR      & \pending{---} & \pending{---} \\
\hline
Fine-tuned base + NN  & \pending{---} & \pending{---} \\
\hline
Fine-tuned base + LR  & \pending{---} & \pending{---} \\
\hline
Fine-tuned tiny + LR  & \pending{---} & \pending{---} \\
\hline
Fine-tuned tiny + NN  & \pending{---} & \pending{---} \\
\hline
\end{tabular}
\caption{Paired bootstrap differences relative to the handcrafted baseline (2{,}000 resamples,
$N=148$, 95\% intervals); negative values indicate performance below handcrafted. Both methods are
evaluated on the same resampled subjects at each iteration. An interval containing zero indicates no
statistically detectable difference. Fine-tuned rows use the last-1-block configuration, the
best-performing depth (Table~\ref{tab:depth}), so these comparisons are the most favorable case for
fine-tuning.}
\label{tab:diff}
\end{table}
\todo{Fill from the final cell of \texttt{04\_bootstrap\_ci.ipynb}, which prints these rows in
LaTeX form ready to paste.}

\begin{figure}[h]
\centering
\includegraphics[width=\textwidth]{bootstrap_comparison.png}
\caption{AD vs.\ control: all seven methods with 95\% bootstrap confidence intervals for accuracy
(left) and AUC (right), 5-fold CV, $N=148$; fine-tune rows use the last-1 block. Red = handcrafted,
green = frozen BrainOmni, blue = fine-tuned BrainOmni. The intervals overlap across all methods.}
\label{fig:bootstrap}
\end{figure}
\todo{Upload \texttt{results\_withseed/bootstrap\_comparison.png} to the Overleaf project (or adjust
the path). Same for the learning-curve figures below.}

\subsection{The result is robust to classifier and dimensionality}
Across $\sim$23 classifiers, logistic regression was the best or tied-best readout of the frozen
features; no classifier robustly outperformed it, and the few that peaked higher did so within
bootstrap noise and only at specific dimensionalities. Accuracy was flat across PCA settings
(1--50 components as well as full dimensionality), confirming the choice of a linear readout and
dimensionality is not driving the result.

\subsection{BrainOmni is data-limited on this task}
\label{sec:datalimited}
Two lines of evidence indicate the bottleneck is data, not optimization. First, the fine-tuning
learning curves (Figure~\ref{fig:curves}) show clear overfitting: training loss falls toward its
floor while the held-out and inner-validation losses \emph{rise} after roughly epoch 25--50, and the
effect worsens with fine-tuning depth (the last-3-block configuration overfits fastest). The
reported accuracies use the best inner-validation epoch, so they are not the overfit endpoint; the
curves are diagnostic. Second, a sample-size learning curve (accuracy vs.\ number of training
subjects) shows the handcrafted baseline saturating early while frozen BrainOmni is still rising at
the full cohort size, i.e.\ on the steep, data-hungry portion of its curve rather than at a plateau.
Together these suggest that 148 subjects is too few to adapt the pretrained backbone: fine-tuning
simply memorizes the training fold.

\begin{figure}[h]
\centering
\includegraphics[width=\textwidth]{nn_curves_cv_base.png}
\caption{5-fold CV learning curves (mean $\pm$1 std across folds), NN head, BrainOmni-base. Top:
accuracy; bottom: loss. Train (blue), inner validation used for epoch selection (orange), held-out
test fold logged for diagnosis only (green). Training loss falls toward its floor while held-out and
validation losses rise, overfitting that worsens from last-1 (left) to last-3 (right).}
\label{fig:curves}
\end{figure}
\todo{Optionally add the sample-size learning curve (\texttt{learning\_curve\_samplesize.png}) as a
second panel; it shows frozen BrainOmni still rising at $N=118$ while handcrafted saturates.}

\section{Discussion}
On this 148-subject cohort, a general-purpose MEG foundation model shows no advantage over six
handcrafted spectral/connectivity features for AD classification, whether used frozen or fine-tuned,
and remains data-limited. Two explanations are consistent with these observations and cannot be
fully separated here: the pretrained representation may not encode AD-specific structure
(a representation gap, since pretraining was not disease-specific), or 148 subjects may be too few
to surface or adapt it (a data limit). The still-rising learning curve favors the latter, while the
persistent tie even at full data is consistent with the former. Either way, the promise of strong
AD detection with minimal task data is not realized here.
\paragraph{Limitations.} A single cohort ($N=148$) and one dataset limit generalizability, and the
wide confidence intervals reflect that sample size. Our preprocessing applied no explicit band-pass
or notch filtering beyond the resample's anti-aliasing low-pass, so low-frequency drift and 50/60\,Hz
line noise were not removed; whether this affects the foundation-model features is a robustness
question worth checking. \todo{Confirm the filtering point with Sri (was the eLife data pre-filtered?)
and, if feasible, report a with/without-filter robustness check for the frozen features.}
\todo{Integrate the MEG-GPT half (Sid): two foundation models tying with / underperforming handcrafted
would strengthen the negative result across architectures.}

\section*{Reproducibility}
All experiments use a fixed seed (123) and a numbered, self-contained notebook pipeline
(\texttt{notebooks\_withseed/00--07}), from frozen-feature extraction through fine-tuning, bootstrap
inference, and robustness sweeps. \todo{Add repository link.}

\end{document}
